\documentclass[11pt,reqno]{amsart}

\usepackage[T1]{fontenc}
\usepackage{lmodern}
\usepackage[expansion=false]{microtype}
\usepackage[a4paper,textwidth=150mm,textheight=242mm,centering,
  headheight=14pt,headsep=7mm,footskip=11mm]{geometry}
\usepackage{amsmath,amssymb,amsthm,mathtools}
\usepackage{hyperref}
\hypersetup{colorlinks=true,linkcolor=blue,citecolor=blue,urlcolor=blue,
  pdftitle={Wehrl-type entropy problem for compact connected semisimple Lie groups},
  pdfauthor={},
  pdfsubject={Haar measure formulation, real Husimi moments, and direct entropy variation}}
\usepackage{bookmark}
\usepackage{braket}
\allowdisplaybreaks[2]
\emergencystretch=2em
\numberwithin{equation}{section}

\theoremstyle{plain}
\newtheorem{theorem}{Theorem}[section]
\newtheorem{proposition}[theorem]{Proposition}
\newtheorem{corollary}[theorem]{Corollary}
\newtheorem{lemma}[theorem]{Lemma}
\theoremstyle{remark}
\newtheorem{remark}[theorem]{Remark}

\newcommand{\C}{\mathbb C}
\newcommand{\N}{\mathbb N}
\newcommand{\E}{\mathbb E}
\newcommand{\Sym}{\operatorname{Sym}}
\newcommand{\Rea}{\operatorname{Re}}
\newcommand{\Ima}{\operatorname{Im}}
\newcommand{\R}{\mathbb R}
\newcommand{\PP}{\mathbb P}
\newcommand{\gfrak}{\mathfrak g}
\newcommand{\rh}{\varrho}
\newcommand{\ii}{\mathrm i}
\newcommand{\dd}{\,\mathrm d}
\newcommand{\Id}{\mathrm I}
\newcommand{\tr}{\operatorname{Tr}}
\newcommand{\Ad}{\operatorname{Ad}}
\newcommand{\conv}{\operatorname{conv}}
\newcommand{\norm}[1]{\lVert#1\rVert}
\newcommand{\ip}[2]{\langle#1,#2\rangle}
\newcommand{\proj}[1]{\lvert#1\rangle\langle#1\rvert}
\newcommand{\SW}{S_{\mathrm W}}

\title[Wehrl-type entropy for compact semisimple groups]{Wehrl-type entropy problem for compact connected semisimple Lie groups}
\date{September 10, 2026}
\subjclass[2020]{Primary 22E46, 43A85; Secondary 81R30, 94A17}
\keywords{Wehrl entropy, coherent states, compact connected semisimple groups,
Casimir operator, Husimi function, highest weights}

\author{Haonan Zhang}
\address{Department of Mathematics, University of South Carolina, Columbia, SC 29208, USA}
\email{haonanzhangmath@gmail.com}
\date{\today}

\begin{document}
\raggedbottom
\begin{abstract}
This paper solves the Wehrl-type entropy problem for arbitrary compact connected semisimple Lie groups.
Let $G$ be a compact connected semisimple Lie group, and let $\pi:G\to U(V_\lambda)$ be a
finite-dimensional irreducible unitary representation associated with the highest weight $\lambda$.
We prove that coherent projectors are the unique minimizers of the Wehrl entropy
over all density matrices on $V_\lambda$. They also uniquely maximize every Husimi power moment of order $p>1$.
The proof uses a second-variation at the extremizer by perturbation in directions associated with Killing fields, similar to the strategy of Frank and Lieb used in \cite{FrankLieb}. Then the problem reduces to the extreme-moment property of highest-weight vectors.

\end{abstract}
\maketitle

\section{Introduction}\label{sec:intro}

Wehrl \cite{Wehrl,Wehrl79} defined a classical analog of the von Neumann entropy
$-\tr(\sigma\log\sigma)$ of density operators using coherent states. For a density operator $\sigma$ on $L^2(\R^n)$, consider the Husimi function
\begin{gather*}
 Q_\sigma(q,p)=\langle q,p|\sigma|q,p\rangle,
\end{gather*}
where $\ket{q,p}\in L^2(\R^n), q,p\in\R^n$ are the Glauber coherent states \cite{Glauber},
\[
 \ket{q,p}(x)=(\pi\hbar)^{-n/4}
 \exp\left(-\frac{|x-q|^2}{2\hbar}
       +\frac{\ii}{\hbar}p\cdot\left(x-\frac q2\right)\right),\qquad x\in \R^n.
       \]
Here $\hbar>0$ is the reduced Planck constant, and we may choose units in which $\hbar=1$. The Husimi function is a probability density on phase space $(\R^{2n},\mathrm d q\,\mathrm d p/(2\pi\hbar)^n)$, and the Wehrl entropy is defined as
\[
 \SW(\sigma)=-\int_{\R^{2n}}Q_\sigma\log Q_\sigma\,
                  \frac{\mathrm d q\,\mathrm d p}{(2\pi\hbar)^n}.
\]
Throughout, we use $0\log0=0$ and take inner products to be conjugate-linear
in the first argument and linear in the second.
Wehrl conjectured \cite{Wehrl79} that coherent
projectors minimize this entropy, and this was solved by Lieb \cite{Lieb}
using sharp Hausdorff--Young and Young convolution inequalities. In fact, he proved the analogous results
 for power functions $x^p, p>1$, and then derived the entropy result by differentiating at $p=1$.
Carlen \cite{Carlen} gave another proof using logarithmic Sobolev inequalities, and his proof works for power functions and $x\log x$.
He also obtained the equality condition for the first time.
Luo \cite{Luo} gave an alternative proof of the original Wehrl entropy inequality based on hypercontractivity.

The above problem has a group-theoretic formulation, and coherent states can be defined
for more general Lie groups \cite{Perelomov}. Wehrl's original problem concerns the Heisenberg group, and it is natural to ask if the same extremal phenomenon occurs for more Lie groups. In particular, Lieb \cite{Lieb} conjectured the same result for the Bloch
coherent states associated with $SU(2)$. After a number of partial results \cite{Bodmann,Gnutzmann,Schupp}, this was solved by Lieb and Solovej
\cite{LS2} by determining the minimum output entropy of a class of quantum channels and a semiclassical limit argument.
Their proof applies to general convex functions, including $x^p$ for $p>1$ and $x\log x$.
The same paper also established the general convex-function inequality for Glauber coherent states by a limiting argument.
Later, the same authors \cite{LS} extended the results to symmetric representations of $SU(N)$.

Lieb and Solovej \cite{LS11} also studied the Wehrl-type entropy problem
for the holomorphic discrete series of $SU(1,1)$; see also \cite{Bandyopadhyay,Bayart}. The corresponding
Wehrl-type inequality was proved for general convex functions by Kulikov
\cite{Kulikov} using the hyperbolic isoperimetric inequality for level
sets of weighted holomorphic functions.
Kulikov also characterized equality for strictly convex functions.
In related work, Nicola and Tilli \cite{NT} proved the Faber--Krahn inequality for the short-time Fourier
transform with Gaussian window.
Frank \cite{Frank} and, independently,
Kulikov--Nicola--Ortega-Cerd\`a--Tilli \cite{KNOT}
gave alternative proofs and characterized the pure-state optimizers for
continuous nonaffine convex functions in the Heisenberg, $SU(2)$, and $SU(1,1)$
settings, when the sharp bound is finite.
Quantitative stability was established for Heisenberg coherent states by
Frank--Nicola--Tilli \cite{FNT}, and equality and stability for symmetric
$SU(N)$ coherent states by Nicola--Riccardi--Tilli \cite{NRT};
see also their alternative stability proof \cite{NRTsecond}.
We refer to the references in these works for more relevant results.

For arbitrary compact connected semisimple groups, Sugita \cite{Sugita,SugitaMoments} proved the
sharp inequalities for integer powers $p\ge2$ by tensor-product methods.
See also van Haastrecht and Zhang \cite[Appendix~A]{HZ} for a detailed
Casimir proof of the Cartan-power equality characterization.

\medskip

In this paper, we provide a solution to the Wehrl-type entropy problem for
arbitrary compact connected semisimple Lie groups. To state the results, let us recall the formulation of the problem first.

Let $G$ be a compact connected semisimple Lie group, and let
$\pi:G\to U(V_\lambda)$ be a finite-dimensional irreducible unitary
representation corresponding to the highest weight $\lambda$.  Write $d=\dim V_\lambda$,
fix a unit highest-weight vector $v_\lambda\in V_\lambda$, and let $\mathrm d g$ be
normalized Haar measure on $G$ so that $\int_G\mathrm d g=1$.
For $g\in G$, set $ v_g=\pi(g)v_\lambda,$ and consider the coherent projectors
\[
 P_g=\proj{v_g}=\pi(g)\proj{v_\lambda}\pi(g)^*.
\]
They satisfy the resolution of identity
\[
 \int_G P_g\dd g=\frac1d\Id.
 \]
 Then for any density matrix $\sigma$ on $V_\lambda$, the Husimi function
 \begin{equation}\label{eq:husimi}
 Q_\sigma(g)=\tr(\sigma P_g)=\ip{v_g}{\sigma v_g}
\end{equation}
is a probability density function on $(G, d\cdot \dd g)$. We define the Wehrl-type entropy by
\begin{equation}
 \SW(\sigma):=-d\int_G Q_\sigma(g)\log Q_\sigma(g)\dd g.
\end{equation}
When $\sigma$ is a pure state $\proj v$,
we write $Q_v=Q_{\proj v}$ and $\SW(v)=\SW(\proj v)$, so that $Q_v(g)=|\ip{v_g}{v}|^2$. The Wehrl-type entropy conjecture asserts that
the minimum of $\SW(\sigma)$ over all density matrices $\sigma$ on $V_\lambda$ is attained at a coherent state $\proj{v_g}$ for some (and for all) $g\in G$.
The generalized problem concerns functions beyond $-x\log x$, as well as the equality conditions.

Our main results are as follows.

\begin{theorem}\label{thm:main}
Let $\Phi:[0,1]\to\R$ be either $\Phi(x)=x^p$ for a fixed real $p>1$,
or $\Phi(x)=x\log x$. For every unit vector $v\in V_\lambda$,
\begin{equation}\label{eq:functional}
 \int_G\Phi(|\langle v_g,v\rangle|^2)\dd g
 \le \int_G\Phi(|\langle v_g,v_\lambda\rangle|^2)\dd g.
\end{equation}
For either choice of $\Phi$, equality holds if and only if
$v=e^{\ii\theta}v_h$ for some $\theta\in\R$ and $h\in G$.
\end{theorem}

It is well-known that when $\Phi$ is convex, the inequality for
general density matrices can be reduced to pure states,
and the same reduction works for the equality condition if $\Phi$ is strictly convex.
See for example \cite[Proof of Corollary~14]{Frank}.

\begin{corollary}\label{cor:density}
For either function $\Phi$ in Theorem~\ref{thm:main} and every density
matrix $\sigma$ on $V_\lambda$,
\begin{equation}\label{eq:densityfunctional}
 \int_G\Phi(Q_\sigma(g))\dd g
 \le \int_G\Phi(Q_{v_\lambda}(g))\dd g.
\end{equation}
Equality holds if and only if $\sigma=P_h$ for some $h\in G$.
In particular,
\begin{equation}\label{eq:entropy}
 \SW(\sigma)\ge\SW(\proj{v_\lambda}).
\end{equation}
\end{corollary}

The idea of the proof is to analyze an optimizer under perturbations in
directions associated with Killing fields.
Suppose that $w_*$ is a unit vector maximizing
$\int_G\Phi(Q_v(g))\dd g$, with value $M$, for $\Phi(x)=x^p$, $p>1$.
Let $T_j=-\ii\,\mathrm d\pi(X_j)$, where $(X_j)_j$ is an orthonormal
basis for the negative Killing form on $\gfrak$. We use
\[
 \gamma_j(t)=\frac{w_*+tT_jw_*}{\norm{w_*+tT_jw_*}},\qquad |t|\text{ small}.
\]
Homogeneity reduces the problem to the scalar functions
\[
 \phi_j(t)=\int_G|\langle v_g,w_*+tT_jw_*\rangle|^{2p}\dd g
              -M\norm{w_*+tT_jw_*}^{2p}.
\]
Using some identities for Casimir operators, we obtain
\[
 \sum_j\phi_j''(0)=4p(p-1)M
 \left(\|\lambda\|^2-\sum_j\langle w_*,T_jw_*\rangle^2\right)\ge 0.
\]
The final inequality follows from the extreme-moment bound in Lemma~\ref{lem:moment}. The choice of $w_*$ also implies
$\phi_j''(0)\le0$ for each $j$, so
$\|\lambda\|^2=\sum_j\langle w_*,T_jw_*\rangle^2$.
This equality holds only if $w_*$ is coherent.

\medskip

The proof strategy is similar to the one used by Frank and Lieb in
\cite{FrankLieb} to prove sharp Hardy--Littlewood--Sobolev inequalities
on the Heisenberg group: starting from an optimizer, they examine second
variations in a similar way, then sum the resulting stability inequalities
and combine them with a spectral comparison to force the optimizer to
have the desired form. The author thanks Rupert Frank for pointing this out.

\medskip

Section~\ref{sec:prelim} recalls the Lie-theoretic preliminaries,
the extreme-moment property, and the Casimir identities.
Section~\ref{sec:variations} collects some technical computations in the second variation argument.
In Section~\ref{sec:proof} we prove the main results.
Some technical details are deferred to Appendices~\ref{app:regularity} and~\ref{app:representation}.

\subsection*{Acknowledgments} The author is partially supported by NSF grant DMS-2453408 and acknowledges the use of OpenAI’s GPT-5.6 and GPT-6. He is grateful to Rupert Frank for valuable comments on an earlier version of the paper, and to Paata Ivanisvili for helpful discussion. The author thanks Jan Philip Solovej for hosting him during a visit to Copenhagen in 2022 and for helpful discussion. Part of the work was finished during a visit to Texas, and he is grateful to Tao Mei and Ping Zhong for kind hospitality.

\section{Preliminaries on Lie groups and Lie algebras}\label{sec:prelim}

For background on highest-weight representations, see \cite{Humphreys,Knapp}.
For the trivial representation, $d=1$ and $Q_\sigma(g)=1$ for every
$g\in G$, so every moment is one and the entropy is zero.
Theorem~\ref{thm:main} is immediate in this case.  We assume $\lambda\ne0$ below.

\subsection{Basic notations}
Let $G$ be a compact semisimple connected Lie group and $\gfrak$ its Lie algebra.
Consider the inner product defined by the negative Killing form
\begin{equation}\label{eq:killingB}
       B(X,Y)=-\operatorname{Kill}(X,Y)
 =-\operatorname{Tr}
       (\operatorname{ad}X\operatorname{ad}Y).
\end{equation}
Since $\gfrak$ is compact and semisimple, $B$ is positive definite and
$\Ad(G)$-invariant.  We use this Killing-form normalization on every
simple factor.  For $X\in\gfrak$, define the self-adjoint operator
\[
 T(X)=-\ii\,\mathrm d\pi(X).
\]
Choose a $B$-orthonormal basis $X_1,\ldots,X_s$ of $\gfrak$ and consider the self-adjoint operators
\[
 T_j=T(X_j),\qquad 1\le j\le s.
\]
Write $h_j(t)=\exp_G(-tX_j)$, $t\in\R$, for the one-parameter subgroup
of $G$ generated by $-X_j$; it acts as
$\pi(h_j(t))=e^{-\ii tT_j}$.

Let $\mathfrak t\subset\gfrak$ be a Cartan subalgebra.
We use real weights $\nu\in\mathfrak t^*$, with the convention
\[
 T(H)w=\nu(H)w,\qquad H\in\mathfrak t,\quad w\in V_\lambda[\nu].
\]
Roots use the same real convention in the adjoint representation.
The inner product $B$ identifies $\mathfrak t^*$ with $\mathfrak t$;
we use the same symbol for a weight and its $B$-dual vector.
With a choice of positive roots, write
\[
 \rh=\frac12\sum_{\alpha>0}\alpha.
\]
Parentheses denote the induced inner product on real weights.

According to Schur's lemma, one has
\begin{equation}\label{eq:normalization}
 \int_G P_g\dd g=\frac1d\Id,\qquad
 \int_G Q_\sigma(g)\dd g=\frac1d.
\end{equation}

\subsection{The extreme-moment property}
The moment of a density matrix $\sigma$ is
\[
 m(\sigma)=\sum_{j=1}^s\tr(\sigma T_j)X_j.
\]
This definition is independent of the orthonormal basis and is
equivariant under $G$; it is the moment map given by expectation
values in our real Lie-algebra convention \cite[Section~2.2]{BCV}.
For a unit vector $v$, write $m(v)=m(\proj v)$, so that
\begin{equation}\label{eq:moment}
 \norm{m(v)}^2=\sum_j\ip v{T_jv}^2.
\end{equation}
The highest-weight projector is
fixed by the torus with Lie algebra $\mathfrak t$, so its moment belongs to $\mathfrak t$.
Its Cartan coordinates are those of the highest weight, giving
$m(v_\lambda)=\lambda$.

For pure states, the following lemma is the minimum-uncertainty
characterization of Delbourgo and Fox \cite{DelbourgoFox}; see also
\cite[Lemma~2]{Sugita}.
The lemma is a consequence of the standard highest-weight
estimates \cite[Theorem~5.5(b),(e), pp.~279--280]{Knapp}; we include
a proof for completeness.

\begin{lemma}[Extreme moments]\label{lem:moment}
For every density matrix $\sigma$,
\begin{equation}\label{eq:momentbound}
 \norm{m(\sigma)}^2\le\norm{\lambda}^2.
\end{equation}
Equality holds exactly for coherent projectors.
 In particular,
 $\norm{m(v)}^2=\norm{\lambda}^2$ if and only if
 $v=e^{\ii\theta}\pi(h)v_\lambda$ for some $\theta\in\R$ and $h\in G$.
\end{lemma}

\begin{proof}
Decompose $V_\lambda=\bigoplus_\nu V_\lambda[\nu]$ into its orthogonal weight
spaces, and let $\Pi_\nu$ be the projection onto $V_\lambda[\nu]$.
Writing $W$ for the Weyl group, the standard weight facts give
$\norm\nu^2\le\norm{\lambda}^2$, with equality exactly for $\nu\in W\lambda$.
For every $\nu\in W\lambda$, the weight space $V_\lambda[\nu]$ is
one-dimensional.

By equivariance, we may conjugate $\sigma$ so that $m(\sigma)\in\mathfrak t$.
Choose an orthonormal basis $H_1,\ldots,H_r$ of $\mathfrak t$.
Since the moment definition is independent of the orthonormal basis,
the coordinates of this conjugated moment are $\tr(\sigma T(H_\ell))$.

 Then
\[
 T(H_\ell)=\sum_\nu\nu(H_\ell)\Pi_\nu,\qquad 1\le\ell \le r,
\]
and we may write
\begin{align*}
 m(\sigma)
 &=\sum_{\ell=1}^r\tr(\sigma T(H_\ell))H_\ell\\
 &=\sum_{\ell=1}^r\sum_\nu
       \tr(\sigma\Pi_\nu)\nu(H_\ell)H_\ell\\
 &=\sum_\nu\tr(\sigma\Pi_\nu)
       \left(\sum_{\ell=1}^r\nu(H_\ell)H_\ell\right)
  =\sum_\nu\tr(\sigma\Pi_\nu)\nu.
\end{align*}
The last equality uses the identification of weights with vectors.
Set $p_\nu=\tr(\sigma\Pi_\nu)\ge0$; since
$\sum_\nu\Pi_\nu=\Id$, we have $\sum_\nu p_\nu=\tr\sigma=1$.
Thus $m(\sigma)$ is the convex combination
\[
 m(\sigma)=\sum_\nu p_\nu\nu,
\]
of weights, and consequently
\[
 \norm{m(\sigma)}^2
 \le\sum_\nu p_\nu\norm\nu^2\le\norm{\lambda}^2.
\]
Equality in the first inequality forces all weights with $p_\nu>0$
to equal one $\nu_0$, by strict convexity; equality in the second gives
$\nu_0\in W\lambda$.  Since $\sigma\ge0$ and
$\tr(\sigma(\Id-\Pi_{\nu_0}))=0$, the range of $\sigma^{1/2}$, and hence
of $\sigma$, lies in $V_\lambda[\nu_0]$.  Thus $\sigma$ is the projector onto
that coherent line.  Conversely, equivariance and
$m(v_\lambda)=\lambda$ give equality for every coherent projector.
\end{proof}

\subsection{Casimir identities}

We collect some identities of Casimir operators that will be used later.
We refer to \cite{Sugita} for related discussion.

\begin{lemma}\label{lem:casimir}
For every $g\in G$ and every $v\in V_\lambda$, we have
\begin{equation}\label{eq:casimir}
 \sum_j T_j^2=(\lambda,\lambda+2\rh)\Id,
\end{equation}
and
\begin{equation}\label{eq:square}
 \sum_j\ip{v_g}{T_jv}^{\,2}=\norm{\lambda}^2\,\ip{v_g}{v}^{\,2}.
\end{equation}
The squares in \eqref{eq:square} are complex squares, not modulus squares.
\end{lemma}
\begin{proof}
For an irreducible representation of highest weight $\nu$, the
standard Casimir eigenvalue formula for a $B$-orthonormal compact
basis is (\cite[Proposition~5.28(b), p.~295]{Knapp})
\[
 -\sum_j\mathrm d\pi_\nu(X_j)^2
       =(\nu,\nu+2\rh)\Id.
\]
 Applying it at $\nu=\lambda$ gives
\eqref{eq:casimir}, since $T_j=-\ii\,\mathrm d\pi(X_j)$.

The highest-weight space of weight $2\lambda$ in $V_\lambda^{\otimes2}$
is one-dimensional.  Complete reducibility therefore gives a unique
irreducible summand $V_{2\lambda}$ containing $v_\lambda^{\otimes2}$.
For each $g\in G$, the square $v_g\otimes v_g$ belongs to this
summand by $G$-invariance.  Set
$\Omega=\sum_jT_j\otimes T_j$.  Then \eqref{eq:casimir} gives
\[
 \sum_j(T_j\otimes\Id+\Id\otimes T_j)^2=2(\lambda,\lambda+2\rh)\Id+2\Omega.
\]
On $V_{2\lambda}$, the Casimir eigenvalue formula gives
$(2\lambda,2\lambda+2\rh)\Id$ for this operator, so
$\Omega(v_g\otimes v_g)=\norm{\lambda}^2 v_g\otimes v_g$.  Self-adjointness now gives
\[
 \sum_j\ip{v_g}{T_jv}^{\,2}
 =\ip{v_g\otimes v_g}{\Omega(v\otimes v)}
 =\norm{\lambda}^2\,\ip{v_g}{v}^{\,2}. \qedhere
\]
\end{proof}

\section{Second variations under perturbations}\label{sec:variations}

Fix $v\in V_\lambda$ and we shall use the perturbation $v+tT_jv$. For this, we write
\[
 A=\ip{v_g}{v},\qquad B_j=\ip{v_g}{T_jv},\qquad C_j=\ip{v_g}{T_j^2v}.
\]
Note that
 $|A+tB_j|^2=|A|^2+2t\Re(\overline A B_j)+t^2|B_j|^2$.

\begin{lemma}\label{lem:scalarvariation}
Suppose $\Phi\in C([0,\infty))\cap C^2((0,\infty))$ and
$z\mapsto\Phi(|z|^2)$ is $C^2$ on $\C$. For $A\ne0$,
\begin{align}\label{eq:affineindividual}
 \left.\frac{\mathrm d^2}{\mathrm dt^2}\right|_0\Phi(|A+tB_j|^2)
 &=2\Phi'(|A|^2)|B_j|^2
   +4\Phi''(|A|^2)\bigl(\Re(\overline A B_j)\bigr)^2.
\end{align}
The right-hand side extends continuously across $A=0$. Consequently,
\begin{align}\label{eq:generalvariation}
 &\sum_j\left.\frac{\mathrm d^2}{\mathrm dt^2}\right|_0
 \int_G\Phi(|\ip{v_g}{v+tT_jv}|^2)\dd g\notag\\
 &\qquad=\int_G\left[4\norm{\lambda}^2|A|^4\Phi''(|A|^2)
 +2(\lambda,\lambda+2\rh)|A|^2\Phi'(|A|^2)\right]\dd g.
\end{align}
\end{lemma}
\begin{proof}
The first formula is the ordinary chain rule.
The $C^2$ hypothesis on the radial function implies that
$\Phi'(x)$ has a finite limit and $x\Phi''(x)\to0$ as $x\downarrow0$,
by comparing its radial and tangential second derivatives.
This justifies the asserted extension.

For the integrated variation, the invariance of Haar measure gives
\[
 U_j(t):=\int_G\Phi(|\ip{v_g}{e^{\ii tT_j}v}|^2)\dd g
        =\int_G\Phi(|A|^2)\dd g.
\]
Indeed, $\ip{v_g}{e^{\ii tT_j}v}=\ip{v_{h_j(t)g}}v$. So $U_j''(0)=0$.
Since
\[
 \ip{v_g}{e^{\ii tT_j}v}
 =A+\ii tB_j-\frac{t^2}{2}C_j+o(t^2),
\]
\[
 |\ip{v_g}{e^{\ii tT_j}v}|^2
 =|A|^2-2t\Im(\overline A B_j)
 +t^2\bigl(|B_j|^2-\Re(\overline A C_j)\bigr)+o(t^2),
\]
the chain rule gives
\[
 0=U_j''(0)=\int_G\left[
 2\Phi'(|A|^2)\bigl(|B_j|^2-\Re(\overline A C_j)\bigr)
 +4\Phi''(|A|^2)\bigl(\Im(\overline A B_j)\bigr)^2\right]\dd g.
\]
Subtract this from the integral of \eqref{eq:affineindividual} and
the $|B_j|^2$ terms cancel. Using
$2(\Re z)^2=|z|^2+\Re(z^2)$, or equivalently
$(\Re z)^2-(\Im z)^2=\Re(z^2)$, the remaining integrand is
\[
 4\Phi''(|A|^2)\Re(\overline A^{\,2}B_j^2)
 +2\Phi'(|A|^2)\Re(\overline A C_j).
\]
This difference extends continuously by zero at $A=0$.
Summing over $j$ and using
\[
 \sum_jB_j^2=\norm{\lambda}^2 A^2,\qquad
 \sum_jC_j=(\lambda,\lambda+2\rh)A
\]
from Lemma~\ref{lem:casimir} gives \eqref{eq:generalvariation} after integration.
All differentiations under the integral are justified by compactness
and the $C^2$ hypothesis.
\end{proof}

\begin{remark}\label{rem:representation}
The unitary second variation is used only for cancellation in the proof above.
It can be avoided by summing and integrating \eqref{eq:affineindividual},
then using the Casimir identities and the weighted identity \eqref{eq:weighted}.
Appendix~\ref{app:representation} states this identity and gives two proofs.
\end{remark}

For normalization, put $r_j(t)=\norm{v+tT_jv}^2$.
\begin{lemma}\label{lem:normvariation}
Let $v$ be a unit vector and $\Phi$ be of class $C^2$ near $1$.
For each $j$,
\begin{equation}\label{eq:normindividual}
 \left.\frac{\mathrm d^2}{\mathrm dt^2}\right|_0\Phi(r_j(t))
 =2\norm{T_jv}^2\Phi'(1)+4\ip v{T_jv}^2\Phi''(1).
\end{equation}
Consequently,
\begin{equation}\label{eq:normcomposition}
 \sum_j\left.\frac{\mathrm d^2}{\mathrm dt^2}\right|_0\Phi(r_j(t))
 =2(\lambda,\lambda+2\rh)\Phi'(1)+4\norm{m(v)}^2\Phi''(1).
\end{equation}
\end{lemma}
\begin{proof}
Expanding $r_j(t)=1+2t\ip v{T_jv}+t^2\norm{T_jv}^2$ gives
$r_j(0)=1$, $r_j'(0)=2\ip v{T_jv}$, and
$r_j''(0)=2\norm{T_jv}^2$. The chain rule proves the first identity.
Summing and using
\[
 \sum_j\norm{T_jv}^2=(\lambda,\lambda+2\rh),\qquad
 \sum_j\ip v{T_jv}^2=\norm{m(v)}^2
\]
proves the second, by Lemma~\ref{lem:casimir} and the definition of $m(v)$.
\end{proof}

\section{Proof of the main results}\label{sec:proof}

We prove Theorem~\ref{thm:main} for two classes of functions separately.

\subsection{The power function \texorpdfstring{$x^p$}{x\textasciicircum p}}\label{sec:powers}

Fix a real $p>1$ and define
\begin{equation}\label{eq:IN}
 I_p(w)=\int_G |\ip{v_g}{w}|^{2p}\dd g,\qquad
 N_p(w)=\norm w^{2p}.
\end{equation}
These functions are homogeneous of degree $2p$, and $I_p$ is
$G$-invariant and phase-invariant.

For fixed $p>1$, a vector $v$, and an index $j$, use $r_j$ from
Section~\ref{sec:variations} and set
\begin{equation}\label{eq:scalarcurves}
 H_j(t)=I_p(v+tT_jv),\qquad n_j(t)=r_j(t)^p.
\end{equation}

\begin{lemma}\label{lem:powerregularity}
The functions $H_j$ and $t\mapsto I_p(e^{\ii tT_j}v)$ are $C^2$.
Their first two derivatives can be taken under the integral.
\end{lemma}
\begin{proof}
The function $g_p(z)=|z|^{2p}$ on $\C\simeq\R^2$ is $C^2$ for
$p>1$, including at zero: its gradient and Hessian are bounded in norm
by $2p|z|^{2p-1}$ and $2p(2p-1)|z|^{2p-2}$, respectively, and
both derivatives extend by zero there.
For a $C^2$ complex-valued function $A(t)$ with $A(t)\ne0$,
\[
 \bigl(|A(t)|^{2p}\bigr)''
 =2p|A|^{2p-2}\bigl(|A'|^2+\Re(\overline A A'')\bigr)
 +4p(p-1)|A|^{2p-4}
                  \bigl(\Re(\overline A A')\bigr)^2.
\]
The second term is bounded in absolute value by
$4p(p-1)|A|^{2p-2}|A'|^2$ and extends continuously by zero at $A=0$.
Take $A(t)=\ip{v_g}{v+tT_jv}$ or $\ip{v_g}{e^{\ii tT_j}v}$.
For small $|t|$, $A,A',A''$ are bounded uniformly in the
group variable $g$, since $\norm{v_g}=1$.  The preceding estimates
and the finiteness of Haar measure justify differentiation under
the integral, also at overlap zeros.
\end{proof}

\begin{proposition}\label{prop:powervariation}
For fixed $p>1$ and $v$, the functions in \eqref{eq:scalarcurves} satisfy
\begin{equation}\label{eq:DI}
 \sum_j H_j''(0)
   =\bigl(2p(\lambda,\lambda+2\rh)+4p(p-1)\norm{\lambda}^2\bigr)I_p(v).
\end{equation}
If $\norm v=1$, then
\begin{equation}\label{eq:DN}
 \sum_j n_j''(0)=2p(\lambda,\lambda+2\rh)+4p(p-1)\norm{m(v)}^2.
\end{equation}
\end{proposition}
\begin{proof}
Lemma~\ref{lem:powerregularity} verifies the regularity hypothesis of
Lemma~\ref{lem:scalarvariation} for $\Phi(x)=x^p$.
Substituting $\Phi'(x)=px^{p-1}$ and
$\Phi''(x)=p(p-1)x^{p-2}$ in \eqref{eq:generalvariation} gives
\[
 \sum_jH_j''(0)
 =\bigl(4p(p-1)\norm{\lambda}^2+2p(\lambda,\lambda+2\rh)\bigr)I_p(v).
\]
For unit $v$, apply Lemma~\ref{lem:normvariation} to the same scalar
function. Since $\Phi'(1)=p$ and $\Phi''(1)=p(p-1)$,
\eqref{eq:normcomposition} gives \eqref{eq:DN} directly.
\end{proof}

Now we are ready to prove the main theorem for the power function.

\begin{proof}[Proof of Theorem~\ref{thm:main}: the power case]
By compactness, $I_p$ attains its maximum $M$ at a unit vector $w_*$.
The value $M$ is positive because $\int_G Q_{w_*}(g)\dd g=1/d$.
Use \eqref{eq:scalarcurves} with $v=w_*$.  For each fixed index $j$, set
\[
 \phi_j(t)=H_j(t)-M n_j(t)
       =I_p(w_*+tT_jw_*)-M\norm{w_*+tT_jw_*}^{2p}.
\]
Homogeneity and maximality give $\phi_j(t)\le0$ for every real $t$,
with $\phi_j(0)=0$.  Hence $\phi_j''(0)\le0$.  Summing over $j$ and
using \eqref{eq:DI} and \eqref{eq:DN}, we obtain
\[
 0\ge\sum_j \phi_j''(0)
 =4p(p-1)M\bigl(\norm{\lambda}^2-\norm{m(w_*)}^2\bigr)\ge0.
\]
Lemma~\ref{lem:moment} gives $\norm{m(w_*)}^2=\norm{\lambda}^2$, so
$w_*=e^{\ii\theta}v_h$ for some $\theta\in\R$ and $h\in G$.
All coherent unit vectors have the same value by invariance, and
therefore $M=I_p(v_\lambda)$.  Every equality vector is a maximizer,
so the same argument makes it coherent.  Conversely, every coherent
vector gives equality.
\end{proof}

\subsection{The logarithmic function \texorpdfstring{$x\log x$}{x log x}}\label{sec:shannon}

For $w\ne0$, define
\begin{equation}\label{eq:logfunctional}
 F(w)=\int_G|\ip{v_g}{w}|^2\log|\ip{v_g}{w}|^2\dd g.
\end{equation}
The resolution of identity gives $\int_G|\ip{v_g}{w}|^2\dd g=\norm w^2/d$, hence
\begin{equation}\label{eq:normalizedlog}
 F\!\left(\frac{w}{\norm w}\right)
 =\frac{F(w)}{\norm w^2}-\frac1d\log\norm w^2.
\end{equation}
The function $x\log x$ is not twice differentiable at zero.
Lemma~\ref{lem:regularity} in Appendix~\ref{app:regularity} justifies
the following differentiation by regularization.

\begin{proposition}\label{prop:variation}
For every unit vector $v$,
\begin{equation}\label{eq:logvariation}
 \sum_j\left.\frac{\mathrm d^2}{\mathrm dt^2}\right|_0F(v+tT_jv)
 =2(\lambda,\lambda+2\rh)F(v)+\frac{2(\lambda,\lambda+2\rh)+4\norm{\lambda}^2}{d}.
\end{equation}
\end{proposition}
\begin{proof}
Set $\Phi_\varepsilon(x)=(x+\varepsilon)\log(x+\varepsilon)$ and
$F_\varepsilon(w)=\int_G\Phi_\varepsilon(|\ip{v_g}{w}|^2)\dd g$, where $\varepsilon>0$.
Its scalar derivatives are
\[
 \Phi_\varepsilon'(x)=1+\log(x+\varepsilon),\qquad
 \Phi_\varepsilon''(x)=\frac1{x+\varepsilon}.
\]
Apply Lemma~\ref{lem:scalarvariation} to obtain
\begin{equation}\label{eq:regularizedvariation}
 \sum_j\left.\frac{\mathrm d^2}{\mathrm dt^2}\right|_0F_\varepsilon(v+tT_jv)
 =\int_G\left[\frac{4\norm{\lambda}^2Q_v^2}{Q_v+\varepsilon}
 +2(\lambda,\lambda+2\rh)Q_v\bigl(1+\log(Q_v+\varepsilon)\bigr)\right]\dd g.
\end{equation}
Lemma~\ref{lem:regularity} permits passage to the second derivatives
as $\varepsilon\downarrow0$. The integrands converge uniformly by
\[
 0\le Q_v-\frac{Q_v^2}{Q_v+\varepsilon}\le\varepsilon,
 \qquad
 0\le Q_v\log(Q_v+\varepsilon)-Q_v\log Q_v
 =Q_v\log\!\left(1+\frac{\varepsilon}{Q_v}\right)\le\varepsilon.
\]
The products involving $Q_v$ are interpreted as zero at $Q_v=0$.
Since $\int_GQ_v\dd g=1/d$, the limit is
\[
 2(\lambda,\lambda+2\rh)F(v)
 +\frac1d\bigl(4\norm{\lambda}^2+2(\lambda,\lambda+2\rh)\bigr),
\]
which is \eqref{eq:logvariation}.
\end{proof}

Now we are ready to prove the main theorem for the logarithmic function.
The inequality also follows directly from the power function result by a limit argument.

\begin{proof}[Proof of Theorem~\ref{thm:main}: the logarithmic case]
Continuity of $x\log x$ on $[0,1]$ makes $F$ continuous on the compact
unit sphere. Let $w_*$ be a maximizer and write $M=F(w_*)$.
For each fixed $j$, put
\[
 H_j(t)=F(w_*+tT_jw_*),\qquad r_j(t)=\norm{w_*+tT_jw_*}^2,
\]
and evaluate the functional under this perturbation, with normalization,
\begin{equation}\label{eq:normalizedlogcurve}
 f_j(t)=F\!\left(\frac{w_*+tT_jw_*}{\norm{w_*+tT_jw_*}}\right)
       =\frac{H_j(t)}{r_j(t)}-\frac1d\log r_j(t).
\end{equation}
We shall show that its second variations force
$\norm{m(w_*)}^2=\norm{\lambda}^2$.
By maximality and Lemma~\ref{lem:regularity},
$f_j'(0)=0$ and $f_j''(0)\le0$. Since $H_j(0)=M$ and $r_j(0)=1$,
the first derivative gives
\[
 H_j'(0)=\left(M+\frac1d\right)r_j'(0).
\]
Differentiating \eqref{eq:normalizedlogcurve} a second time and using
this relation yields
\begin{equation}\label{eq:normalizedlogsecond}
 f_j''(0)=H_j''(0)-\left(M+\frac1d\right)r_j''(0)
                   -\frac{(r_j'(0))^2}{d}.
\end{equation}
Summing \eqref{eq:normalizedlogsecond} and using
Lemma~\ref{lem:normvariation} for the two norm terms and
Proposition~\ref{prop:variation} for $\sum_jH_j''(0)$, we obtain
\begin{equation}\label{eq:entropycontact}
 0\ge\sum_jf_j''(0)
   =\frac4d\bigl(\norm{\lambda}^2-\norm{m(w_*)}^2\bigr)\ge0.
\end{equation}
Lemma~\ref{lem:moment} therefore makes $w_*$ coherent.
Haar invariance gives the same value for all coherent unit vectors,
so $M=F(v_\lambda)$. Every equality vector is a maximizer and is
therefore coherent; conversely, every coherent vector gives equality.
This proves \eqref{eq:functional} for $\Phi(x)=x\log x$.
\end{proof}

\appendix
\section{Regularity of the logarithmic variation}\label{app:regularity}

Since the perturbation $w+tT_jw$ has zero second derivative, no
acceleration term occurs in the logarithmic second variation. The only issue is differentiation at overlap zeros.
\begin{lemma}\label{lem:log-regularity}\label{lem:regularity}
Fix $w\ne0$ and $j$, and put $A=\ip{v_g}w$, $B_j=\ip{v_g}{T_jw}$.
The function $H(t)=F(w+tT_jw)$ is twice differentiable at zero, with
\begin{equation}\label{eq:logaffineindividual}
 H''(0)=\int_G\left[2(1+\log|A|^2)|B_j|^2
       +4\frac{(\Re(\overline A B_j))^2}{|A|^2}\right]\dd g.
\end{equation}
Moreover, for $F_\varepsilon$ in the proof of
Proposition~\ref{prop:variation},
$\left.\frac{\mathrm d^2}{\mathrm dt^2}\right|_0F_\varepsilon(w+tT_jw)
\to H''(0)$ as $\varepsilon\downarrow0$.
\end{lemma}
\begin{proof}
The real-analytic function $|A|^2$ is not identically zero, since
$\int_G|A|^2\dd g=\norm w^2/d>0$. Connectedness, compactness, and
\cite[Proposition~3.1]{ShiZhang} imply $\log|A|^2\in L^1(G)$;
its zero set has Haar measure zero.

Set $L(z)=z(1+\log|z|^2)$ for $z\ne0$, and $L(0)=0$.
On each bounded disk,
\begin{equation}\label{eq:logincrement}
 |L(z+h)-L(z)|\le C|h|(1+|\log|z||),\qquad z\ne0.
\end{equation}
If $|h|\le |z|/2$, integrate the ordinary real derivatives of $L$
along the segment, where their size is bounded by
$C(1+|\log|z||)$. Otherwise $|z|<2|h|$ and $|z+h|<3|h|$;
the bound follows from $|L(\zeta)|\le C|\zeta|(1+|\log|\zeta||)$
and boundedness of $s|\log s|$ for $0\le s\le3$ after scaling by $|h|$.

Since $z\mapsto|z|^2\log|z|^2$ is $C^1$, compactness gives
\[
 H'(t)=2\Re\int_G\overline{L(A+tB_j)}B_j\dd g.
\]
By \eqref{eq:logincrement}, the difference quotient of this integrand
is bounded by $C|B_j|^2(1+|\log|A||)$, an integrable function.
Dominated convergence therefore gives \eqref{eq:logaffineindividual}.
We set the integrand to zero on the null set $\{A=0\}$.

For the regularized derivative, replace $\log|A|^2$ by
$\log(|A|^2+\varepsilon)$ and the denominator by
$|A|^2+\varepsilon$. For $0<\varepsilon\le1$, use
\[
 \frac{(\Re(\overline A B_j))^2}{|A|^2+\varepsilon}\le |B_j|^2,
 \qquad |1+\log(|A|^2+\varepsilon)|\le C+|\log|A|^2|.
\]
Dominated convergence proves the asserted limit.
\end{proof}

\section{Weighted identities and two proofs}\label{app:representation}

For $v\in V_\lambda$, use the notation $A=\ip{v_g}{v}$ and
$B_j=\ip{v_g}{T_jv}$ from Section~\ref{sec:variations}.
Under the hypotheses of Lemma~\ref{lem:scalarvariation}, the weighted identity is
\begin{align}\label{eq:weighted}
 &\int_G\bigl(\Phi'(|A|^2)+|A|^2\Phi''(|A|^2)\bigr)
                    \sum_j|B_j|^2\dd g\notag\\
 &\qquad=\norm{\lambda}^2\int_G|A|^4\Phi''(|A|^2)\dd g
      +(\lambda,\lambda+2\rh)\int_G|A|^2\Phi'(|A|^2)\dd g.
\end{align}
The integrands are interpreted by continuous extension at $A=0$.
Taking $\Phi(x)=x^p$ and dividing by $p$ gives, for every real $p>1$,
\begin{equation}\label{eq:weightedpower}
 p\int_G|A|^{2p-2}\sum_j|B_j|^2\dd g
       =(\lambda,p\lambda+2\rh)\int_G|A|^{2p}\dd g.
\end{equation}
We give two proofs of these identities below.

\subsection{Representation theory}

The Hermitian generators of
$\overline\pi$ are $-\overline{T_j}$, so
\begin{align*}
 \sum_j|B_j|^2
 &=\left\langle v_g\otimes\overline{v_g},
   \left(\sum_jT_j\otimes\overline{T_j}\right)
   (v\otimes\overline v)\right\rangle,\\
 \sum_jT_j\otimes\overline{T_j}
 &=(\lambda,\lambda+2\rh)\Id
   -\frac12\sum_j
      (T_j\otimes\Id-\Id\otimes\overline{T_j})^2.
\end{align*}
The squared operator sum is the Casimir of
$\pi\otimes\overline\pi$; on each irreducible summand with highest
weight $\nu$, it acts by $(\nu,\nu+2\rh)$.

The pointwise sum generally involves several irreducible components.
Its weighted integral can nevertheless be evaluated using ordinary
integer tensor powers. We use the following
formulas  in \cite[Section~III, equations~(14) and~(23)]{Sugita},
with our normalization of Haar measure. Assume first that $\norm v=1$. For each integer
$n\ge1$, let $P_n$ be the orthogonal projection onto the Cartan
component $V_{n\lambda}\subset V_\lambda^{\otimes n}$ generated by
$v_\lambda^{\otimes n}$, and put
\[
 T_j^{(n)}=\sum_{\ell=1}^n
 \Id^{\otimes(\ell-1)}\otimes T_j\otimes\Id^{\otimes(n-\ell)}.
\]
The projection $P_n$ commutes with these total generators, and
\[
 \ip{v_g^{\otimes n}}{T_j^{(n)}v^{\otimes n}}=nA^{n-1}B_j,\qquad
 \int_G\proj{v_g^{\otimes n}}\dd g
   =\frac{P_n}{\dim V_{n\lambda}}.
\]
The second identity is Schur's lemma, and also gives
$I_n(v)=\norm{P_nv^{\otimes n}}^2/\dim V_{n\lambda}$.
Here the same integral definition gives $I_1(v)=\int_G|A|^2\dd g=1/d$.

The Casimir eigenvalue on $V_{n\lambda}$ now yields
\begin{align*}
 n^2\int_G|A|^{2n-2}\sum_j|B_j|^2\dd g
 &=\frac{\sum_j\norm{T_j^{(n)}P_nv^{\otimes n}}^2}
         {\dim V_{n\lambda}}\\
 &=\frac{(n\lambda,n\lambda+2\rh)\norm{P_nv^{\otimes n}}^2}
         {\dim V_{n\lambda}}\\
 &=(n\lambda,n\lambda+2\rh)I_n(v).
\end{align*}
Thus the desired weighted identity holds for every integer $n\ge1$.
By linearity, these identities imply, for every polynomial $f$,
\begin{align*}
 \int_G f(|A|^2)\sum_j|B_j|^2\dd g
 =\norm{\lambda}^2\int_G|A|^2f(|A|^2)\dd g+2(\lambda,\rh)
       \int_G\left(\int_0^{|A|^2}f(x)\,\mathrm dx\right)\dd g.
\end{align*}
Since $|A|^2\in[0,1]$ and $\sum_j|B_j|^2$ is bounded, both sides
are continuous in the uniform norm of $f$ on $[0,1]$.
Uniform polynomial approximation therefore extends this identity to
every continuous $f$. Taking $f(x)=x^{p-1}$ gives, for every real $p>1$,
\[
 \int_G|A|^{2p-2}\sum_j|B_j|^2\dd g
 =\left(\norm{\lambda}^2+\frac{2(\lambda,\rh)}p\right)I_p(v)
 =\frac{(\lambda,p\lambda+2\rh)}p I_p(v).
\]
For general $\Phi$ satisfying the hypotheses of Lemma~\ref{lem:scalarvariation}, take
$f(x)=\Phi'(x)+x\Phi''(x)$, continuously extended at zero.
Then $\int_0^q f(x)\,\mathrm dx=q\Phi'(q)$, since
$x\Phi'(x)\to0$ at zero. The continuous weighted identity above
is exactly \eqref{eq:weighted}. Rescaling extends it to arbitrary $v$,
by applying the unit-vector identity to $x\mapsto\Phi(\norm v^2x)$;
$v=0$ follows directly. The mixed representation identifies the
pointwise operator, while the Cartan tensor powers evaluate its weighted integral.

\subsection{Integration by parts}
Fix $v$ throughout. Let $\mathcal X_j$ act on scalar functions by
\[
 \mathcal X_j f(g)=\left.\frac{\mathrm d}{\mathrm dt}\right|_0
                         f(\exp_G(tX_j)g),\qquad
 \Delta=\sum_j\mathcal X_j^2.
\]
For the bi-invariant metric given by the negative Killing form,
these vector fields are orthonormal and divergence-free. We have
$\mathcal X_j A=-\ii B_j$ and
$\Delta A=-(\lambda,\lambda+2\rh)A$.
Writing $q=|A|^2$ for this calculation gives
\begin{align*}
 \Delta q&=2\sum_j|B_j|^2-2(\lambda,\lambda+2\rh)q,\\
 |\nabla q|^2&=4\sum_j\bigl(\Im(\overline A B_j)\bigr)^2
       =2\left(q\sum_j|B_j|^2-\norm{\lambda}^2q^2\right).
\end{align*}
Integration by parts on the compact group without boundary gives
\[
 0=\int_G\Delta\Phi(q)\dd g
   =\int_G\left[\Phi'(q)\Delta q+\Phi''(q)|\nabla q|^2\right]\dd g.
\]
Substituting the two scalar identities and dividing by $2$ proves
\eqref{eq:weighted}. At zeros, $\Phi(q)$ is $C^2$ by hypothesis,
and the chain-rule expression extends continuously. In particular,
for $\Phi(x)=x^p$, $p>1$, one may directly use
$0=\int_G\Delta(|A|^{2p})\dd g$.

\newcommand{\etalchar}[1]{$^{#1}$}
\begin{thebibliography}{KNOCT25}
\raggedright
\setlength{\baselineskip}{10pt}
\enlargethispage{2\baselineskip}

\bibitem[Ban09]{Bandyopadhyay}
J.~Bandyopadhyay,
Optimal concentration for $SU(1,1)$ coherent state transforms and an
analogue of the Lieb--Wehrl conjecture for $SU(1,1)$,
\emph{Comm. Math. Phys.} \textbf{285} (2009), 1065--1086.
\href{https://doi.org/10.1007/s00220-008-0618-y}{doi:10.1007/s00220-008-0618-y};
\href{https://arxiv.org/abs/0711.0013}{arXiv:0711.0013}.

\bibitem[BBH{\etalchar{+}}19]{Bayart}
F.~Bayart, O.~F.~Brevig, A.~Haimi, J.~Ortega-Cerd\`a, and K.-M.~Perfekt,
Contractive inequalities for Bergman spaces and multiplicative Hankel forms,
\emph{Trans. Amer. Math. Soc.} \textbf{371} (2019), 681--707.
\href{https://doi.org/10.1090/tran/7290}{doi:10.1090/tran/7290}.

\bibitem[BCV21]{BCV}
A.~Botero, M.~Christandl, and P.~Vrana,
Large deviation principle for moment map estimation,
\emph{Electron. J. Probab.} \textbf{26} (2021), article 79, 1--23.
\href{https://doi.org/10.1214/21-EJP636}{doi:10.1214/21-EJP636}.

\bibitem[Bod04]{Bodmann}
B.~G.~Bodmann,
A lower bound for the Wehrl entropy of quantum spin with sharp high-spin
asymptotics,
\emph{Comm. Math. Phys.} \textbf{250} (2004), 287--300.
\href{https://doi.org/10.1007/s00220-004-1146-z}{doi:10.1007/s00220-004-1146-z}.

\bibitem[Car91]{Carlen}
E.~A.~Carlen,
Some integral identities and inequalities for entire functions and their
application to the coherent state transform,
\emph{J. Funct. Anal.} \textbf{97} (1991), 231--249.
\href{https://doi.org/10.1016/0022-1236(91)90022-W}{doi:10.1016/0022-1236(91)90022-W}.

\bibitem[DF77]{DelbourgoFox}
R.~Delbourgo and J.~R.~Fox,
Maximum weight vectors possess minimal uncertainty,
\emph{J. Phys. A: Math. Gen.} \textbf{10} (1977), L233--L235.
\href{https://doi.org/10.1088/0305-4470/10/12/004}{doi:10.1088/0305-4470/10/12/004}.

\bibitem[FL12]{FrankLieb}
R.~L.~Frank and E.~H.~Lieb,
Sharp constants in several inequalities on the Heisenberg group,
\emph{Ann. of Math. (2)} \textbf{176} (2012), no.~1, 349--381.
\href{https://doi.org/10.4007/annals.2012.176.1.6}
{doi:10.4007/annals.2012.176.1.6}.

\bibitem[FNT25]{FNT}
R.~L.~Frank, F.~Nicola, and P.~Tilli,
The generalized Wehrl entropy bound in quantitative form,
\emph{J. Eur. Math. Soc.}, published online 2025.
\href{https://doi.org/10.4171/JEMS/1674}{doi:10.4171/JEMS/1674}.

\bibitem[Fra23]{Frank}
R.~L.~Frank,
Sharp inequalities for coherent states and their optimizers,
\emph{Adv. Nonlinear Stud.} \textbf{23} (2023), article 20220050, 28~pp.
\href{https://doi.org/10.1515/ans-2022-0050}{doi:10.1515/ans-2022-0050}.

\bibitem[Gla63]{Glauber}
R.~J.~Glauber,
Coherent and incoherent states of the radiation field,
\emph{Phys. Rev.} \textbf{131} (1963), 2766--2788.
\href{https://doi.org/10.1103/PhysRev.131.2766}{doi:10.1103/PhysRev.131.2766}.

\bibitem[G{\.Z}01]{Gnutzmann}
S.~Gnutzmann and K.~\.{Z}yczkowski,
R\'enyi--Wehrl entropies as measures of localization in phase space,
\emph{J. Phys. A: Math. Gen.} \textbf{34} (2001), 10123--10139.
\href{https://doi.org/10.1088/0305-4470/34/47/317}{doi:10.1088/0305-4470/34/47/317};
\href{https://arxiv.org/abs/quant-ph/0106016}{arXiv:quant-ph/0106016}.

\bibitem[Hum72]{Humphreys}
J.~E.~Humphreys,
\emph{Introduction to Lie Algebras and Representation Theory},
Graduate Texts in Mathematics, vol.~9, Springer, 1972.
\href{https://doi.org/10.1007/978-1-4612-6398-2}{doi:10.1007/978-1-4612-6398-2}.

\bibitem[Kna02]{Knapp}
A.~W.~Knapp,
\emph{Lie Groups Beyond an Introduction}, second edition,
Progress in Mathematics, vol.~140, Birkh\"auser, 2002.
\href{https://www.math.stonybrook.edu/~aknapp/download/Beyond2-clickable.pdf}
{Author-hosted digital second edition with 2023 errata}.

\bibitem[KNOCT25]{KNOT}
A.~Kulikov, F.~Nicola, J.~Ortega-Cerd\`a, and P.~Tilli,
A monotonicity theorem for subharmonic functions on manifolds,
\emph{Adv. Math.} \textbf{479} (2025), article 110423.
\href{https://doi.org/10.1016/j.aim.2025.110423}{doi:10.1016/j.aim.2025.110423}.

\bibitem[Kul22]{Kulikov}
A.~Kulikov,
Functionals with extrema at reproducing kernels,
\emph{Geom. Funct. Anal.} \textbf{32} (2022), 938--949.
\href{https://doi.org/10.1007/s00039-022-00608-5}{doi:10.1007/s00039-022-00608-5}.

\bibitem[Lie78]{Lieb}
E.~H.~Lieb,
Proof of an entropy conjecture of Wehrl,
\emph{Comm. Math. Phys.} \textbf{62} (1978), 35--41.
\href{https://doi.org/10.1007/BF01940328}{doi:10.1007/BF01940328}.

\bibitem[LS14]{LS2}
E.~H.~Lieb and J.~P.~Solovej,
Proof of an entropy conjecture for Bloch coherent spin states and its
generalizations,
\emph{Acta Math.} \textbf{212} (2014), 379--398.
\href{https://doi.org/10.1007/s11511-014-0113-6}{doi:10.1007/s11511-014-0113-6}.

\bibitem[LS16]{LS}
E.~H.~Lieb and J.~P.~Solovej,
Proof of the Wehrl-type entropy conjecture for symmetric $SU(N)$ coherent
states,
\emph{Comm. Math. Phys.} \textbf{348} (2016), 567--578.
\href{https://doi.org/10.1007/s00220-016-2596-9}{doi:10.1007/s00220-016-2596-9}.

\bibitem[LS21]{LS11}
E.~H.~Lieb and J.~P.~Solovej,
Wehrl-type coherent state entropy inequalities for $SU(1,1)$ and its
$AX+B$ subgroup,
in \emph{Partial Differential Equations, Spectral Theory, and Mathematical
Physics---The Ari Laptev Anniversary Volume},
EMS Series of Congress Reports, EMS Press, Berlin, 2021, 301--314.
\href{https://doi.org/10.4171/ECR/18-1/18}{doi:10.4171/ECR/18-1/18};
\href{https://arxiv.org/abs/1906.00223}{arXiv:1906.00223}.

\bibitem[Luo00]{Luo}
S.~Luo,
A simple proof of Wehrl's conjecture on entropy,
\emph{J. Phys. A: Math. Gen.} \textbf{33} (2000), no.~16, 3093--3096.
\href{https://doi.org/10.1088/0305-4470/33/16/303}{doi:10.1088/0305-4470/33/16/303}.

\bibitem[NRT26a]{NRTsecond}
F.~Nicola, F.~Riccardi, and P.~Tilli,
An elementary approach to Wehrl-type entropy bounds in quantitative form,
\emph{Proc. Lond. Math. Soc.} (3) \textbf{133} (2026), article e70182.
\href{https://doi.org/10.1112/plms.70182}{doi:10.1112/plms.70182}.

\bibitem[NRT26b]{NRT}
F.~Nicola, F.~Riccardi, and P.~Tilli,
The Wehrl-type entropy conjecture for symmetric $SU(N)$ coherent states:
cases of equality and stability,
\emph{Duke Math. J.}, published online 2026, 18~pp.
\href{https://doi.org/10.1215/00127094-2025-0070}{doi:10.1215/00127094-2025-0070};
\href{https://arxiv.org/abs/2412.10940}{arXiv:2412.10940}.

\bibitem[NT22]{NT}
F.~Nicola and P.~Tilli,
The Faber--Krahn inequality for the short-time Fourier transform,
\emph{Invent. Math.} \textbf{230} (2022), no.~1, 1--30.
\href{https://doi.org/10.1007/s00222-022-01119-8}{doi:10.1007/s00222-022-01119-8}.

\bibitem[Per86]{Perelomov}
A.~Perelomov,
\emph{Generalized Coherent States and Their Applications},
Texts and Monographs in Physics, Springer-Verlag, Berlin, 1986.
\href{https://doi.org/10.1007/978-3-642-61629-7}{doi:10.1007/978-3-642-61629-7}.

\bibitem[Sch99]{Schupp}
P.~Schupp,
On Lieb's conjecture for the Wehrl entropy of Bloch coherent states,
\emph{Comm. Math. Phys.} \textbf{207} (1999), 481--493.
\href{https://doi.org/10.1007/s002200050734}{doi:10.1007/s002200050734};
\href{https://arxiv.org/abs/math-ph/9902017}{arXiv:math-ph/9902017}.

\bibitem[Sug02]{Sugita}
A.~Sugita,
Proof of the generalized Lieb--Wehrl conjecture for integer indices larger
than one,
\emph{J. Phys. A: Math. Gen.} \textbf{35} (2002), L621--L626.
\href{https://doi.org/10.1088/0305-4470/35/42/105}{doi:10.1088/0305-4470/35/42/105};
\href{https://arxiv.org/abs/nlin/0208007}{arXiv:nlin/0208007}.

\bibitem[Sug03]{SugitaMoments}
A.~Sugita,
Moments of generalized Husimi distributions and complexity of many-body
quantum states,
\emph{J. Phys. A: Math. Gen.} \textbf{36} (2003), 9081--9103.
\href{https://doi.org/10.1088/0305-4470/36/34/310}{doi:10.1088/0305-4470/36/34/310}.

\bibitem[SZ25]{ShiZhang}
Z.~Shi and R.~Zhang,
Sobolev differentiability properties of logarithmic modulus of real analytic functions,
\emph{J. Funct. Anal.} \textbf{288} (2025), no.~6, article 110804.
\href{https://doi.org/10.1016/j.jfa.2024.110804}{doi:10.1016/j.jfa.2024.110804}.

\bibitem[vHZ26]{HZ}
R.~van Haastrecht and G.~Zhang,
Wehrl inequalities for matrix coefficients of holomorphic discrete series,
\emph{J. Geom. Anal.} \textbf{36} (2026), article 324.
\href{https://doi.org/10.1007/s12220-026-02573-z}{doi:10.1007/s12220-026-02573-z};
\href{https://arxiv.org/abs/2412.18382}{arXiv:2412.18382}.
The compact-group appendix proves the integer-power inequality and the
Cartan-power equality characterization.

\bibitem[Weh78]{Wehrl}
A.~Wehrl,
General properties of entropy,
\emph{Rev. Mod. Phys.} \textbf{50} (1978), 221--260.
\href{https://doi.org/10.1103/RevModPhys.50.221}{doi:10.1103/RevModPhys.50.221}.

\bibitem[Weh79]{Wehrl79}
A.~Wehrl,
On the relation between classical and quantum-mechanical entropy,
\emph{Rep. Math. Phys.} \textbf{16} (1979), 353--358.
\href{https://doi.org/10.1016/0034-4877(79)90070-3}{doi:10.1016/0034-4877(79)90070-3}.

\end{thebibliography}
\end{document}
